\scanpage{46}

\uline{Rmq} Comme premier invariant d'un système $\rho,\rho',\rho''$, il y a \unsure{les|le} \inserted{triplet des} \uline{ordres} $(\unsure{\nu|v},\nu',\nu'')$ de $\rho,\rho',\rho''$. Les op. de $\unsure{\mathfrak S_3|\mathcal S_3}$ sur $H^3$ correspondent \unsure{évidemment|évidemt} aux op. de $\mathfrak S_3$ sur $\mathbb N^3$. Notons $\nu,\nu',\nu''\geq 2$ et

\diagram{$\begin{array}{c@{\ }l}
\begin{array}{c}\mathrm I\\\mathrm{II}\\\mathrm{III}\end{array}&
\left\{\begin{array}{ll}
2\leq\nu,\nu',\nu''\leq3&\text{cas }\unsure{A_4|\mathcal A_4|\mathfrak A_4}\\
2\leq\nu,\nu',\nu''\leq4&\text{cas }\mathfrak S_4\\
2\leq\nu,\nu',\nu''\leq5&\text{cas }A_5
\end{array}\right.
\end{array}$}

On verra que dans le premier cas, \deleted{\unsure{les|le}} triplet contient \deleted{\unsure{contient}} $3$, dans le deuxième il contient $4$, dans le troisième il contient $5$ -- donc le triplet numérique implique déjà dans lequel des cas I, II, III on se trouve (\deleted{\unsure{en reg|en regardant}} suivant que son sup est $3,4$ ou $5$). Dans les cas I II, on voit que $(\rho,\rho',\rho'')$ est déterminé mod iso par son caractère numérique. Dans le cas icosaédr. III, il n'en est plus tout à fait ainsi (pas même pour l'équivalence modulo $\mathfrak S_5$ opérant) -- il faut préciser un peu le critère numérique [en distinguant les deux classes de conjugaison d'él. d'ordre $5$ dans $A_5$ \ldots]

\deleted{3) Cas de $A_4$.}

\deleted{\uline{Théorème} Soit $\Pi$ un tétraèdre \inserted{comb.} orienté, $\unsure{\mathbb H|H}$ le gpe de ses automorphismes \inserted{(orientés)}. À tt sommet $s\in S$ de $\Pi$ est associé un $\rho_s\in H$, él. d'ordre $3$ et on trouve ainsi une bijection de $S$ \deleted{\unsure{dans}} l'ens. des \inserted{$8$} éléments d'ordre $3$ de $H$, soit $H_3^+$ son image, de sorte}

% archive sheet number: 45
